
\textbf{Théorème 4.6.}
\uline{Soient} @@M0@@ \uline{un schéma quasi-compact ayant un faisceau ample}, @@M1@@ \uline{un entier} @@M2@@ \uline{et} @@M3@@ \uline{une immersion régulière de codimension} @@M4@@. \uline{Alors}, \uline{pour tout} @@M5@@, \uline{on a} :
@@M6@@

Remarquons tout d'abord que @@M7@@. En effet, @@M8@@ est dense dans @@M9@@ : s'il n'en était pas ainsi, @@M10@@ contiendrait un ouvert affine, et sur cet ouvert, l'idéal @@M11@@ définissant @@M12@@ serait un nilidéal, donc un idéal localement nilpotent puisqu'il est de type fini. Or ceci est impossible, puisque les @@M13@@ sont des @@M14@@-Modules localement libres de rang @@M15@@. L'augmentation dans @@M16@@ d'un élément de @@M17@@ est une fonction localement constante sur @@M18@@; étant nulle sur @@M19@@, elle est nulle sur @@M20@@ entier.

Pour tout groupe abélien @@M21@@, nous poserons @@M22@@. Soit @@M23@@ ; on peut écrire @@M24@@ sous la forme :
@@M25@@
avec @@M26@@ pour tout @@M27@@, @@M28@@, et @@M29@@. Les @@M30@@ peuvent s'exprimer comme combinaison à coefficients entiers des classes @@M31@@ d'un nombre fini de @@M32@@-Modules localement libres de type fini @@M33@@. Soit @@M34@@ la sous-@@M35@@-@@M36@@-algèbre de @@M37@@ engendrée par les @@M38@@ et par @@M39@@. On a alors :
@@M40@@
Par ailleurs, la filtration de @@M41@@-anneau de @@M42@@ est discrète (VI 6.2); soit @@M43@@ tel que @@M44@@. Enfin, on peut, quitte à décomposer @@M45@@ en une somme finie d'ouverts disjoints et à faire la démonstration dans chaque ouvert (le @@M46@@ commutant aux sommes directes finies), supposer les @@M47@@ de rang constant @@M48@@. L'anneau d'augmentation de @@M49@@ est alors @@M50@@.

